\documentclass[11pt,a4paper]{article}
\usepackage[margin=1in]{geometry}
\usepackage[T1]{fontenc}
\usepackage[activate=false]{microtype}
\usepackage{amsmath,amssymb,amsthm}
\usepackage{enumitem,booktabs,tabularx,array}
\usepackage[hidelinks]{hyperref}
\usepackage{fancyhdr}
\hypersetup{pdftitle={A cubic-communication variant of Multimmit},pdfauthor={}}
\setlength{\parindent}{0pt}
\setlength{\parskip}{5pt plus 1pt}
\setlength{\emergencystretch}{2em}
\setlist{itemsep=3pt,topsep=4pt,parsep=0pt,leftmargin=*}
\newtheorem{theorem}{Theorem}
\newtheorem{lemma}[theorem]{Lemma}
\newtheorem{proposition}[theorem]{Proposition}
\theoremstyle{definition}
\newtheorem{definition}[theorem]{Definition}
\newcommand{\SC}{\mathsf{SC}}
\newcommand{\NC}{\mathsf{NC}}
\newcommand{\Tips}{\mathsf{Tips}}
\newcommand{\Ord}{\mathsf{Ord}}
\newcommand{\Emit}{\mathsf{Emit}}
\newcommand{\GST}{\mathrm{GST}}
\newcommand{\enc}{\operatorname{enc}}
\newcommand{\supp}{\mathsf{support}}
\newcommand{\FF}{\mathbb{F}}
\pagestyle{fancy}
\fancyhf{}
\fancyhead[L]{\small Multimmit: cubic communication}
\fancyhead[R]{\small Construction and proof}
\fancyfoot[C]{\thepage}
\setlength{\headheight}{14pt}

\begin{document}
\begin{center}
{\LARGE\bfseries A cubic-communication variant of Multimmit}\par
\vspace{6pt}
{\large Small view-exit certificates and recovery from existing votes}
\end{center}

\textbf{Outcome.} The construction below gives worst-case
\(O(Kn^2+n^3)\) consensus communication for \emph{each logical view}, hence
\(O(n^3)\) for \(K=\Theta(n)\). It works for every configured fault bound
\(n\ge5f+1\), including \(f\ll n\). It changes the view-progression rules,
retains full atomic votes and the original extraction rules, and supplies an
explicit mechanism for recovering every referenced V-QC. The conservative
timeout is \(4\Delta\). A correct leader's block finalises within \(5\delta\)
of the first correct entry into its view, and within \(2\delta\) of the
leader's proposal transmission.

\textbf{Roadmap.} Small support certificates first restore sequential view
progress without forwarding full V-QCs. A background rule then resolves each
departed view into a genuine V-QC or a nullification. The remaining obstacle
is recovery of an exact, possibly selectively disclosed, parent V-QC. Every
correct recipient eventually knows all selected correct votes already, so
only at most \(f\) records can be missing. We encode the repair information
for those records, and spread it through the correct validators whose votes
establish the proposal's validity. Sections~\ref{sec:recovery}--\ref{sec:live}
give the reconstruction and liveness proofs; Section~\ref{sec:cost} counts
every new send.

\textbf{Scope and status.} This is a proposed replacement for the
communication modification in the reviewed supplement, with a mathematical
proof for the fixed-committee core. It is not a claim that the payment/epoch
extension or the paper's separate lower bounds have been repaired. The
quorum and extraction arguments of the core are reused under the explicitly
proved witness property below. The companion manuscript version is
\texttt{multimmit-31\_10\_26\_1300.tex}~\cite{multimmit}; this
document replaces the earlier dissemination supplement in full. The construction has been checked
against the counterexamples in the third correctness review; the coding
component has also been tested independently. It is not a mechanised proof.

\section{Model and the unchanged core interface}
\label{sec:model}

There are \(n\) validators, at most \(f\) Byzantine, and \(K\) transaction
chains. The committee is fixed and \(n\ge5f+1\). After global stabilisation,
a message sent by a correct validator at time \(s\) reaches every intended
correct recipient by \(\max\{s,\GST\}+\delta\), with known
\(\Delta\ge\delta\). This is the delivery convention in the revised
manuscript. Local enabled actions are exhausted within each timeslot.
Cryptographic verification and local coding are treated as local computation,
as are the paper's other certificate computations.

Hashes are collision resistant. Ordinary signatures, aggregate signatures
and the existing \((2f+1)\)-threshold-signature scheme retain the paper's
security assumptions. New threshold messages use a separate domain tag; no
new threshold key or succinct-proof assumption is needed.

The following core objects and rules are retained.
\begin{itemize}
\item A \emph{full vote} for a leader block \(\ell\) contains the complete
  position and extension vectors, authenticated atomically by one ordinary
  signature. With fixed proposal and extension windows, its message body
  has size \(O(K)\). A correct validator sends at most one such vote per
  view and never both a vote and a novote for that view.
\item A \emph{full V-QC} contains \(n-f\) authenticated vote/novote records
  from distinct validators, including at least \(2f+1\) full votes for its
  designated leader block. It has worst-case size \(O(Kn)\).
  An \emph{L-QC} contains \(n-f\) full votes for one leader block.
\item A \emph{nullification} \(\NC_u\) is the existing
  \((2f+1)\)-threshold signature on \((\mathsf{nullify},u)\).
  An unvoted validator's timeout sends both a novote and a nullify share.
  A validator that has voted for \(\ell\) sends a nullify share only after
  obtaining \(2f+1\) messages from distinct validators, each a novote,
  nullify share, or full vote for a different leader block of that view.
\item A directly valid proposal for view \(v\) references a valid full
  V-QC \(Q\) from a view \(q<v\), proposes chains well formed relative to
  \(\Tips(Q)\), and is accompanied locally by nullifications for every
  \(q<u<v\). The leader-signature and no-known-equivocation checks remain.
\item \(\Tips(Q)\) is computed from the full votes in \(Q\), and
  \(\Ord(Q)\) recursively uses the exact parent QCs named by the leader
  blocks. The rules deriving finalised tips and emitted prefixes from full
  vote pools are unchanged. No small certificate below is an input to these
  extraction functions.
\end{itemize}

Retain the original signed messages, recovered QCs, manifests and recovery
shards needed by later references, as in the core's retained-message model.
An independent storage or pruning optimisation needs a separate argument.

Communication counts prescribed transmissions by correct validators.
Arbitrary Byzantine flooding cannot be bounded. As in the manuscript, hash,
signature and bounded metadata sizes are suppressed. Transaction bodies and
their chain-layer dissemination remain outside this consensus-traffic bound.
The \(O(n^3)\) claim is total communication, not a per-validator bound.

\section{Small certificates for sequential progress}
\label{sec:progress}

Every full vote broadcast now includes the complete \(O(K)\)-size leader
block \(\ell\) and a threshold share on
\[
  (\supp,v,H(\ell)).
\]
The support share is sent \emph{only together with that validator's actual
full vote for the same block}. It is not sent merely because a certificate
was received. The additional commitment in \(\ell\), defined in
Section~\ref{sec:recovery}, is included in \(H(\ell)\) and therefore in both
the ordinary vote and this support share.

\begin{definition}[Support certificate]
A support certificate \(\SC_v(\ell)\) consists of the leader block \(\ell\)
for view \(v\) and a valid \((2f+1)\)-threshold signature on
\((\supp,v,H(\ell))\). Its size is \(O(K)\).
\end{definition}

The modified rules are as follows.
\begin{enumerate}[label=\textbf{P\arabic*.}]
\item \textbf{Forward small exit evidence.} Each correct validator forwards
  its first valid support certificate for each view once to all validators.
  It also forwards nullifications as in the original protocol. Support
  certificates for the same view need not designate the same block;
  uniqueness is not assumed. There is no all-to-all full-V-QC forwarding.
\item \textbf{Traverse views sequentially.} A validator exits its current
  view \(v\) only on \(\SC_v(\ell)\) or \(\NC_v\), and enters \(v+1\).
  Later-view evidence is cached. Several consecutive views may be processed
  within one timeslot, but no intermediate view is skipped. Remove the
  advancement-certificate and capped status/pull rules of the supplement.
\item \textbf{Vote before a support exit.} On \(\SC_v(\ell)\), an unvoted,
  non-nullified validator first casts its full vote for \(\ell\), using its
  own position and extension vectors, and broadcasts that vote, \(\ell\),
  and its support share. It then exits. The one-vote and vote/novote guards
  remain. A nullification exit need not cast a vote.
\item \textbf{Full V-QCs do not independently trigger exit or bypass voting.}
  They remain valid anchors and extraction objects. Direct proposal-validity
  voting and support-certificate voting are the two voting routes.
\item \textbf{Continue old-view nullification.} Retain each entered view's
  voted block and nullified flag after leaving it. Continue applying the
  original opposition-triggered nullify rule to these old views, sending
  at most one nullify share per view. Old timers need not continue.
\item \textbf{Use a \(4\Delta\) active-view timeout.} Direct-validity voting
  and support-certificate vote-before-exit have priority over timeout at
  the same timeslot. The timeout sends a novote and nullify share only when
  the original unvoted/non-nullified guards hold.
\item \textbf{Leaders supply genuine anchors.} A correct leader proposes
  once, as soon as it has a valid full V-QC anchor and the required gap
  nullifications. It chooses the greatest available anchor view below its
  current view, with a fixed tie-break, and broadcasts the full anchor with
  its proposal. Its one full-anchor broadcast is charged to this proposal
  view. It can compute the recovery manifest locally from the anchor.
\end{enumerate}

The ordinary signed vote remains intact: the support share certifies a
separate view-progression assertion. A support certificate may exist before
any \(n-f\)-message V-QC can be assembled. The next leader therefore waits
for a real anchor. The proof in Section~\ref{sec:live} shows that this wait
terminates without a large-object pull protocol.

Attaching \(\ell\) to every full-vote broadcast is necessary. Otherwise,
recipients might collect enough matching votes without obtaining their
proposal body: a voter may already have used its single support-certificate
forwarding allowance for a different block. This attachment costs only
\(O(K)\) per vote.

\section{Exact-QC recovery from existing votes}
\label{sec:recovery}

The main observation is that all correct full votes and novotes are already
broadcast. For any particular valid V-QC, every correct recipient eventually
obtains every selected correct signer's record. Thus at most \(f\) selected
records can remain missing. It is unnecessary to redisperse \(n\) records
to repair those missing \(f\).

\subsection{Canonical representation and a precise manifest}

Fix a canonical, injective QC representation. A full QC \(Q\) is identified
by its designated leader block, its selected signer bitmap, the selected
\emph{complete unsigned message records} in signer-index order, and its
aggregate signature. Define \(H(Q)\) on that canonical representation.
The current standard-vote/deviation compression may still be used on the
wire, provided decompression yields this representation. Different lossless
wire compressions of one QC then have one identifier.

Individual constituent signatures are not part of the recovery rows: the
aggregate representation may not expose them. The aggregate signature is
retained and verified against all reconstructed message records.

Write \(X_i\) for the unsigned record selected from signer \(i\), with a
zero row in an unselected slot. Records include their message type and all
signed vote content; use an unambiguous fixed-width padding convention.
Choose a finite field large enough for the codes below and represent each
row by a vector in \(\FF^B\). Here \(B\) represents \(O(K)\) message data,
up to field-symbol padding. The selected-signers bitmap distinguishes absent
slots from actual records.

The \emph{manifest} \(M(Q)\) consists of:
\begin{enumerate}
\item \(H(Q)\), the designated leader block, the selected-signers bitmap,
  the aggregate signature, and the canonical format/length parameters;
\item the \(n\) indexed row hashes \(H(X_1),\ldots,H(X_n)\);
\item the \(n\) indexed recovery-shard hashes defined below.
\end{enumerate}
All hashes bind their context and index. The \(O(K+n)\)-size manifest uses
full hash vectors, avoiding Merkle proofs within the cubic budget.

\subsection{Two levels of erasure coding}

Assume first that \(f>0\). Fix an \(f\times n\) Cauchy matrix \(A\) over
\(\FF\), so every square submatrix is invertible. For example, choose
pairwise distinct field elements \(a_1,\ldots,a_f,b_1,\ldots,b_n\) and set
\(A_{ji}=1/(a_j-b_i)\). Compute the \(f\) parity rows
\begin{equation}
  R_j=\sum_{i=1}^{n} A_{ji}X_i,\qquad 1\le j\le f.
  \label{eq:inner}
\end{equation}
These \(f\) rows recover any \(f\) missing original rows when the other
rows are known. The invertibility property follows from the Cauchy
determinant formula; it is the standard Cauchy Reed--Solomon
construction~\cite{plank}.

Concatenate the parity rows into \(R\), of length \(fB\) field symbols.
Pad and split \(R\) into \(f+1\) equal vectors, then apply a fixed
\((n,f+1)\) MDS erasure code, giving indexed shards
\(c_1(Q),\ldots,c_n(Q)\). Any \(f+1\) distinct shards reconstruct \(R\).
For example, each coordinate of the input vectors can be treated as the
coefficients of a degree-at-most-\(f\) polynomial, evaluated at \(n\)
distinct field points; \(f+1\) evaluations recover the coefficients by
interpolation~\cite{rao}. Each shard has length
\begin{equation}
  \left\lceil\frac{fB}{f+1}\right\rceil=O(B)=O(K),
  \label{eq:shard}
\end{equation}
with the small symbol-padding term understood. This bound remains valid
when \(f\) is much smaller than \(n\). For \(f=0\), set every \(c_i(Q)\)
to the empty string and use its indexed hash in the manifest. Omit both
coding steps, but still broadcast the manifest: every selected record is
already broadcast by a correct signer. Genesis is globally known and
requires no recovery.

\subsection{Bind and transmit the recovery information}

Replace a leader block's tuple by
\[
  \ell=(v,q,r,C),\qquad q=H(Q),\quad r=H(M(Q)).
\]
Both the leader signature and \(H(\ell)\) bind \(r\). A \emph{direct-valid
vote} is a vote made after checking the full proposal-validity conditions,
including possession of the exact full \(Q\) and recomputation of
\(M(Q)\), its row hashes, parities, and shard hashes. A correct direct-valid
voter \(i\) broadcasts
\begin{equation}
  (v,H(\ell),M(Q),i,c_i(Q))
  \label{eq:repair}
\end{equation}
before or atomically with its ordinary vote and support share. The proposal
body accompanies the vote. A recovery packet may be combined with that
same broadcast. A validator sends at most one recovery packet in a view,
for its one direct-valid vote; it does not send responses to arbitrary
requests. A support-route vote without the checked full anchor sends no
such packet. Previously initiated sends continue after view advancement.

Recipients check \(H(M(Q))=r\), the shard index and its committed hash,
and retain only hash-matching candidate row records. They can use the
correct voters' original broadcasts as these records. When enough matching
shards and records arrive, they reconstruct \(Q\), check \(H(Q)=q\), and
verify its full V-QC structure and aggregate signature. They also check the
deterministic manifest on the reconstructed QC. A wrong Byzantine record
becomes an erasure; it is never used as an unchecked equation input.

\begin{lemma}[Direct-valid witnesses]\label{lem:witness}
If \(\ell\) has a support certificate, or is designated by a valid full
V-QC or L-QC, then at least \(f+1\) correct validators have cast
direct-valid votes for \(\ell\).
\end{lemma}
\begin{proof}
Work in any finite execution prefix in which the hypothesised certificate
already exists. If a correct support-route vote for \(\ell\) exists in
that prefix, consider the earliest
such vote in causal order. Its support certificate requires \(2f+1\)
earlier support shares, including at least \(f+1\) from correct validators.
Each of those shares accompanies an actual earlier full vote for \(\ell\).
By minimality, those correct votes are direct-valid votes. If there is no
correct support-route vote for \(\ell\), then the at least \(f+1\) correct
votes guaranteed by the hypothesised certificate are themselves
direct-valid. This argument uses the rule that observing a support
certificate alone does not cause a fresh support share without a full vote.
\end{proof}

\begin{lemma}[Exact recovery]\label{lem:recovery}
For every proposal \(\ell\) covered by Lemma~\ref{lem:witness}, every
correct validator eventually obtains its exact referenced full V-QC \(Q\).
If a valid support certificate, full V-QC or L-QC for \(\ell\) is held by
a correct validator at time \(s\), recovery completes everywhere by
\(\max\{s,\GST\}+\delta\).
\end{lemma}
\begin{proof}
Apply Lemma~\ref{lem:witness} to the execution prefix ending at \(s\).
Its \(f+1\) correct direct-valid witnesses hold the same exact \(Q\) and
recompute the same deterministic manifest. Their distinct indexed recovery
shards were broadcast by \(s\). All correct recipients therefore obtain
\(f+1\) authenticated shards and recover \(R\) by the stated deadline.

Every selected correct row of \(Q\) was also broadcast by its signer before
\(Q\) could be checked by those witnesses. These rows arrive by the same
deadline. Unselected zero rows are known from the bitmap. The remaining
unknown selected rows form a set \(U\) of size \(u\le f\). Subtract known
contributions from the first \(u\) parity equations to obtain
\[
  R_j-\sum_{i\notin U}A_{ji}X_i
       =\sum_{i\in U}A_{ji}X_i,\qquad 1\le j\le u.
\]
The coefficient matrix is an invertible Cauchy submatrix. Solve
coordinatewise to recover all missing rows, then reconstruct the exact
canonical QC using its aggregate signature and header from the manifest.
Hash and full certificate checks authenticate the result. When \(f=0\),
all rows arrive directly and the same conclusion holds without parity.
\end{proof}

Applying this argument to the finite descending chain of referenced views
recovers every QC required to compute \(\Ord\). Different Byzantine vote
variants cause no ambiguity: the manifest selects the exact row hashes
used by the named QC. Multiple proposals in one faulty-leader view do not
multiply any correct sender's allowance; each correct validator votes and
sends its recovery contribution at most once in that view.

\section{Safety argument and proof changes}
\label{sec:safety}

The extraction certificates and their complete atomic vote contents are
unchanged. The new manifest hash is merely an additional field of the
leader-block identifier. All position and extension computations still use
the same chain proposal \(C\) and the voter's own DA history.

The essential proof adaptation is Lemma~\ref{lem:witness}. In the original
leader-chain and availability arguments, a V-QC-triggered vote is traced
back to earlier correct voters that checked the proposal. The new support
route has the same required \(f+1\) direct-valid witnesses. Consequently,
every certified proposal inherits an actually checked full parent QC,
well-formed chain proposal, and gap nullifications. A support certificate
never supplies positions or substitutes for a full QC in a quorum count.

One-vote-per-view and vote/novote exclusivity follow from the retained
historical flags. If a block receives an L-QC, a support certificate for a
different block of that view is impossible: at most \(f\) correct validators
can lie outside the L-QC's correct supporters, and at most \(f\) faulty
validators can add conflicting shares, fewer than \(2f+1\).
The original no-nullification argument also remains
valid with background processing: if a block obtains the original large
vote quorum, consider the first correct supporter that might later
nullify. Its \(2f+1\) opposing messages cannot exist, using the same
correct-signer count as in the core proof. This argument concerns all
messages ever sent for that view, so leaving the view does not change it.

The leader-chain argument can therefore be rerun with the witness lemma.
The safe-extension and quorum-intersection calculations continue to use
the original full votes and full QCs; their thresholds do not change.
The emission-prefix argument uses the current pool's settledness, as
identified in the third review: reaching the relevant second-pass slot
without halting gives
\[
  \beta_j+(n-|P|)\le f,
\]
and completing the first pass gives \(k_j^*=m_j\). These establish
settledness relative to \(P\), allowing Safe extension~(b) to supply the
required tip equality. This is the independent local correction to retain
when invoking consistency.

Thus the operational core's consistency and inclusion arguments transfer
with this stated change to the causal voting lemma and the existing local
pool correction. The separate universal lower-bound proofs are not used
to establish the communication construction.

\section{View progress, usable anchors and finalisation}
\label{sec:live}

\begin{lemma}[Sequential synchronisation]\label{lem:sync}
If the first correct entry into view \(v\) occurs at time \(t\), every
correct validator enters a view at least \(v\) by
\(\max\{t,\GST\}+\delta\).
\end{lemma}
\begin{proof}
The entering validator has processed every view \(u<v\). For each it has
forwarded a support certificate or nullification by time \(t\). The
forwarded support certificate need not be the same variant it eventually
used: either is adequate exit evidence. By the stated deadline every
correct validator has exit evidence for every \(u<v\), and processes the
views sequentially within the timeslot, casting any prescribed votes
before advancing.
\end{proof}

\begin{lemma}[Resolution of a departed view]\label{lem:resolve}
Let \(s\) be the first correct entry into view \(u+1\), and put
\(S=\max\{s,\GST\}\). By \(S+3\delta\), every correct validator has
either a valid full V-QC for \(u\) or a nullification for \(u\).
\end{lemma}
\begin{proof}
By Lemma~\ref{lem:sync}, every correct validator has passed through \(u\)
by \(S+\delta\). If any correct validator holds a nullification for \(u\)
by \(S+2\delta\), forwarding puts it everywhere by \(S+3\delta\).

Otherwise no exit by \(S+\delta\) can have used a nullification. Each
correct validator exited on a support certificate. It had already voted,
cast its prescribed vote before exit, or had previously sent a timeout
novote. A validator that nullified without voting could only have done so
by that timeout rule. Hence every correct validator broadcast a full vote
or novote for \(u\) by \(S+\delta\). Everyone obtains all these messages,
with the full proposal bodies attached to votes, by \(S+2\delta\).

If one proposal has at least \(2f+1\) correct votes, choose those and enough
other correct messages to obtain a valid \(n-f\)-message V-QC. Everyone
can assemble one locally. Otherwise each proposal has at most \(2f\)
correct votes. Every correct voter then has at least
\[
  (n-f)-2f=n-3f\ge2f+1
\]
correct messages not supporting its chosen block. The retained background
rule makes it send its nullify share, if it has not already done so.
Novoters already sent theirs. These shares produce a nullification
everywhere by \(S+3\delta\).
\end{proof}

The deadline split in this proof matters. A nullification exit can occur
without a vote or novote. The proof only counts all correct view messages
after excluding that case; it does not silently assume that every departure
produces a vote.

\begin{lemma}[A correct leader obtains an anchor]\label{lem:anchor}
If the first correct entry into view \(v\) is at \(t\ge\GST\), a correct
leader can issue a valid proposal by \(t+3\delta\).
\end{lemma}
\begin{proof}
Every first correct exit from an earlier view occurs by \(t\). Apply
Lemma~\ref{lem:resolve} simultaneously to every such view. By
\(t+3\delta\), the leader has a full V-QC or nullification for each
\(u<v\). Select the greatest view \(q<v\) for which it holds a full
V-QC, using genesis if necessary. Every \(q<u<v\) must then have a
nullification. Thus a valid anchor and all required gap evidence are
available. If it could propose earlier, rule P7 already instructed it to do
so. If it has not yet proposed, it is still in the view: no support
certificate can exist without earlier correct votes for its signed
proposal, and no nullification can precede the \(4\Delta\) timeout without
a first correct nullifier or a competing correct vote. Thus P7 can execute.
Coding the manifest introduces local computation but no additional
communication prerequisite.
\end{proof}

\begin{theorem}[Correct-leader finalisation]\label{thm:final}
Let the first correct entry into a correct-leader view \(v\) occur at
\(t\ge\GST\), and let \(\tau\) be the leader's proposal time. Then
\[
  t\le\tau\le t+3\delta,
\]
all correct validators vote by \(\tau+\delta\le t+4\delta\), and all
obtain an L-QC and finalise the leader block by
\(\tau+2\delta\le t+5\delta\).
\end{theorem}
\begin{proof}
The leader signs only one block for the view. No correct validator can vote
for a different block: a first such vote is neither directly valid nor
support-triggered, since the latter would require earlier correct votes
for that same different block. No correct timeout occurs before
\(t+4\Delta\ge t+4\delta\). No earlier correct nullification can occur
either: a first such event would need more opposing messages than the
\(f\) faulty validators can supply, since correct competing votes,
timeout novotes and earlier correct nullifications are absent.

By \(\tau+\delta\), all correct validators have reached at least \(v\)
by Lemma~\ref{lem:sync}. They also have the proposal, its full anchor and
all gap nullifications: the leader has broadcast the anchor and proposal,
and forwarded the nullifications by \(\tau\). The deterministic manifest
check is local. Validators still in \(v\) vote, with priority over any
same-timeslot timeout. Any validator already beyond \(v\) must have
exited on this proposal's support certificate and voted first. Hence all
correct votes are sent by \(\tau+\delta\), and all arrive by
\(\tau+2\delta\), giving the required \(n-f\)-vote pools.

The current anchor is supplied in full. All of its proper reference-chain
ancestors are also recoverable by \(\tau+\delta\): their required correct
recovery packets and selected correct records were broadcast before the
top anchor could exist, by the causal proof of Lemma~\ref{lem:recovery}.
Thus missing ordering metadata does not postpone the conclusion. Delivery
of transaction bodies retains the core's separate data-retrieval convention.
\end{proof}

\begin{lemma}[Progress through any view]\label{lem:exit}
If the first correct entry into view \(v\) occurs at \(t\ge\GST\), all
correct validators leave it by \(t+4\Delta+3\delta\).
\end{lemma}
\begin{proof}
If a correct validator exits by \(t+4\Delta+2\delta\), synchronisation
gives the result. Otherwise every correct validator remains in the view,
having entered by \(t+\delta\). Each sends its full vote or timeout novote
by \(t+4\Delta+\delta\), so everyone has all correct messages by
\(t+4\Delta+2\delta\). If one proposal has \(2f+1\) correct votes, their
support shares produce an exit certificate. Otherwise the opposing-message
argument in Lemma~\ref{lem:resolve} produces a nullification within one
more delay.
\end{proof}

The same argument with \(\max\{t,\GST\}\) starts from views whose first
entry predates GST. Induction gives eventual traversal of every view.
Round-robin leaders therefore yield recurring correct-leader finalisations.
Together with exact reference recovery and the unchanged block-inclusion
argument, this supplies the core's eventual transaction liveness.

\section{Worst-case communication for each view}
\label{sec:cost}

All limits below are permanent per-view limits; they do not reset when a
request arrives or another validator advances. A send relating to an old
view is charged to the view named by its action. A recovery contribution
and the leader's full-anchor broadcast are charged to the view of the
proposal that references that anchor.

\begin{center}
\small
\begin{tabularx}{\textwidth}{@{}Xrr@{}}
\toprule
Transmission class & Copies per view & Size per copy\\
\midrule
Full votes, with proposal bodies and support shares & \(O(n^2)\) & \(O(K)\)\\
Novotes and nullify shares, including late closure & \(O(n^2)\) & \(O(1)\)\\
First support-certificate forwarding & \(O(n^2)\) & \(O(K)\)\\
Nullification forwarding & \(O(n^2)\) & \(O(1)\)\\
Correct leader's one full-anchor proposal & \(O(n)\) & \(O(Kn)\)\\
One recovery contribution per direct-valid voter & \(O(n^2)\) & \(O(K+n)\)\\
\bottomrule
\end{tabularx}
\end{center}

\begin{theorem}[Per-view bound]\label{thm:cost}
For every execution and logical view, correct validators send
\(O(Kn^2+n^3)\) consensus data in \(O(n^2)\) messages. In particular, for
\(K=\Theta(n)\), the total is \(O(n^3)\) per view in the worst case,
including before GST. The bound includes recovery of named ancestor QCs.
\end{theorem}
\begin{proof}
Each correct validator sends at most one ordinary vote or novote, at most one
nullify share, at most one recovery contribution, and at most one broadcast
of each small certificate type per view. There is only one correct leader
for a given view, and it sends at most one full anchor to each recipient.
Sum the table. A Byzantine leader may equivocate, but cannot cause a correct
validator to exceed its one-vote or one-recovery-contribution limit. A
private Byzantine QC does not trigger full-QC forwarding. Recovery causes
no additional request-response traffic.
\end{proof}

This is an independent cap for each logical view, not an amortised argument
that pays for an expensive view with cheap later ones. Sequential catch-up
may process a backlog of many old views in one physical interval; the
theorem does not bound such a physical interval by the cost of one view.
Through \(V\) logical views the total is \(O(V(Kn^2+n^3))\).

For \(K\ll n\), the present construction should be advertised as
\(O(Kn^2+n^3)\), not as \(O(Kn^2)\). The indexed hash manifests contribute
the \(n^3\) term. Restoring security-parameter and field-symbol sizes adds
their conventional factors; field-padding terms of \(O(n^2\log n)\) are
absorbed by \(O(n^3)\) under the paper's fixed cryptographic-size convention.
A nonleader's prescribed outgoing traffic is \(O((K+n)n)\); a correct
leader additionally sends \(O(Kn^2)\) for its full-anchor proposal.

\section{Transferred results}\label{sec:transfer}

The construction is designed so that the manuscript's results
transfer with one substitution. The following proposition makes
the transfer explicit, so that this document, together with the
manuscript's proofs, is checkable line by line.

\begin{proposition}[Carried results]\label{prop:carry}
Under the construction of
Sections~\ref{sec:progress}--\ref{sec:recovery}, the
manuscript's [Designation], [Safe extension], [Availability],
[Block inclusion], [Consistency] and [Emission is a prefix] hold
verbatim, and [Liveness] holds with the constants of
Section~\ref{sec:live}.
\end{proposition}

\begin{proof}
Three kinds of dependence on the dissemination rules occur in the
manuscript's proofs. First, pure counting over signed objects:
full votes, V-QCs, L-QCs, pools and nullifications are
byte-for-byte the manuscript's objects, and no support
certificate enters any extraction, quorum or threshold count
(rule P4), so [Designation], [Safe extension] and every
quorum-intersection step are untouched. Second, causal
valid-proposal inferences: where a proof traces a certificate
back to a correct voter that checked the proposal in full --- the
anchor and base availability steps of [Availability], and the
frontier-endorsement step of [Block inclusion] ---
Lemma~\ref{lem:witness} supplies at least $f+1$ correct
direct-valid voters behind every support certificate, V-QC and
L-QC, each of which checked exactly the original validity
conditions, full-anchor possession included; substituting these
witnesses for the single checked voter the original arguments
extract, the proofs go through unchanged. Third, exact
references: where [Consistency]'s recursive $\Ord$ computation
needs the exact parent QCs named on the reference chain,
Lemma~\ref{lem:recovery} delivers them to every correct
processor, and [Emission is a prefix] uses the pool-settledness
inference in the manuscript's corrected form, reaching a
second-pass slot without halting giving the pool continuation
condition, the completed first pass giving $k^*_j=m_j$, hence
settledness in $P$ and Safe extension~(b). For [Liveness]:
Lemma~\ref{lem:exit} and induction traverse every view,
round-robin leadership supplies recurring correct-leader views,
Theorem~\ref{thm:final} finalises each within $5\delta$ of its
first correct entry after GST, and [Block inclusion]'s
hypotheses are met as above, so the manuscript's transaction-%
liveness argument runs with these lemmas in place of the
replaced synchronisation, exit and correct-leader lemmas.
\end{proof}

\section{Integration, tradeoffs and validation}

The replacement should be specified as one coherent variant. In particular,
do not retain the supplement's skipping rule or its capped QC retrieval
alongside these rules. Also remove the old full-V-QC bypass-vote and exit
routes from the replacement pseudocode; their purpose is served by the
new support certificate. Full V-QCs continue to determine ordering.

\begin{samepage}
The proof changes are concentrated in the following places:
\begin{itemize}
\item extend the leader-block identifier and signature to bind the recovery
  manifest, and define canonical QC hashing;
\item add support shares, the compact certificate, sequential vote-before-exit,
  historical nullification state, and the \(4\Delta\) timeout;
\item replace the causal valid-proposal inference in the leader-chain and
  availability/inclusion proofs by Lemma~\ref{lem:witness};
\item replace the view-synchronisation, view-exit and correct-leader timing
  lemmas with Section~\ref{sec:live};
\item replace the unspecified named-QC retrieval convention with
  Lemma~\ref{lem:recovery}, and retain the pool-settledness proof correction.
\end{itemize}
\end{samepage}

The main cost is additional local coding and hashing, plus the larger
conservative timeout. There is no mandatory extra communication round
between a correct leader's proposal and finalisation. An anchor may require
up to the proved \(3\delta\) preparation interval after first view entry.
The note does not claim the original \(2\Delta\) timeout or original
\(3\delta\) entry-to-finalisation bound for this variant. Simple matrix
implementations of the coding can be computationally expensive; the
communication theorem does not establish an implementation throughput figure.

\textbf{Checks performed.} The view-progress argument was checked against
the earlier alternating skipped-voter execution, conflicting support
certificates, late nullifications, an old view with missing messages after
a nullification exit, and arbitrary \(n\ge5f+1\). Exact recovery was checked
for selectively disclosed Byzantine records, different Byzantine variants,
aggregate-only QC signatures, and \(f\ll n\).

An independent finite-field implementation passed \(11{,}930\)
reconstruction cases: exhaustive missing-row and recovery-shard subsets for
\((n,f)=(6,1)\) and \((11,2)\); deterministic samples for \((16,3)\),
\((50,9)\), and \((37,1)\); and \(f=0\) cases. It also rejected \(247\)
tampered row/shard inputs and passed five combined filtered-recovery cases.
These are coding-component checks, not distributed-protocol model checking.
The proof above supplies the protocol argument.

\begin{thebibliography}{9}
\bibitem{multimmit}
Anonymous author(s).
\newblock Multimmit: Faster finality for ordering and payments.
\newblock Under submission, 2026.
\bibitem{plank} James S. Plank.
\emph{Optimizing Cauchy Reed-Solomon Codes for Fault-Tolerant Storage
Applications}. Technical Report CS-05-569, University of Tennessee, 2005,
Section 2.
\url{https://web.eecs.utk.edu/~jplank/plank/papers/CS-05-569.pdf}.
\bibitem{rao} Anup Rao.
\emph{Lecture 4: Reed-Solomon Codes}. University of Washington, 2019,
pp.~2--3.
\url{https://homes.cs.washington.edu/~anuprao/pubs/codingtheory/lecture4.pdf}.
\end{thebibliography}

\end{document}
